\documentclass[11pt]{article}
\usepackage[a4paper,margin=26mm]{geometry}
\usepackage[T1]{fontenc}
\usepackage{lmodern}
\usepackage{amsmath,amssymb,amsthm}
\usepackage{microtype}
\usepackage{booktabs}
\usepackage{xurl}
\usepackage{hyperref}
\hypersetup{colorlinks=true,linkcolor=blue,urlcolor=blue,citecolor=blue,
 pdftitle={An axial algebra whose automorphism group does not preserve its Frobenius form: an answer to a question of Shpectorov},
 pdfauthor={Horace Dong}}
\newtheorem{theorem}{Theorem}
\newtheorem{lemma}{Lemma}
\theoremstyle{remark}
\newtheorem{remark}{Remark}
\newcommand{\repoversion}{1.20.0}
\newcommand{\F}{\mathbb{F}}
\setlength{\emergencystretch}{1.5em}

\title{An axial algebra whose automorphism group\\ does not preserve its Frobenius form:\\ an answer to a question of Shpectorov}
\author{Horace Dong\\ \small Independent researcher\\
\small\texttt{maxwelldhx@gmail.com}\\
\small ORCID: \href{https://orcid.org/0009-0002-7554-7548}{0009-0002-7554-7548}}
\date{13 November 2026}
\begin{document}
\maketitle

\begin{abstract}
In their survey of open questions on axial algebras, Gorshkov and Shpectorov
record a question of Shpectorov: is there an axial algebra with a Frobenius
form that is not preserved by the full automorphism group of the algebra? They
ask for a non-trivial example, and especially for one whose space of Frobenius
forms is $1$-dimensional. We give a $4$-dimensional example over the field
$\F_{11}$. It is a simple algebra without identity, generated by four primitive
axes for a $C_2$-graded fusion law on $\{1,5,7,10\}$; every bilinear form that
associates with its product is a multiple of one non-degenerate symmetric
form~$\beta$, and an automorphism of order~$2$ multiplies $\beta$ by~$-1$. The
fusion law is not of Jordan or Monster type, and the field has positive
characteristic; the question does not exclude either. Theorem~1 is formally
verified in Lean~4. The example and its proofs were found and formalized by AI
agents working under the author's direction; their role is described at the end
of the paper.
\end{abstract}

\section{Introduction}

We use the definitions of Gorshkov and Shpectorov~\cite[Section~2]{gs}, which
we recall briefly. Let $A$ be a commutative, not necessarily associative algebra
over a field~$\F$. A \emph{fusion law} is a set $\mathcal F\subseteq\F$ with
$1\in\mathcal F$ and a map $\ast$ from $\mathcal F\times\mathcal F$ to subsets
of~$\mathcal F$. For $a\in A$ and $\lambda\in\F$ let $A_\lambda(a)$ be the
$\lambda$-eigenspace of the multiplication map $L_a\colon u\mapsto au$, and for
$\Lambda\subseteq\F$ let $A_\Lambda(a)=\sum_{\lambda\in\Lambda}A_\lambda(a)$. A
non-zero idempotent $a$ is an \emph{$\mathcal F$-axis} if
$A=A_{\mathcal F}(a)$ and $A_\lambda(a)A_\mu(a)\subseteq A_{\lambda\ast\mu}(a)$
for all $\lambda,\mu\in\mathcal F$; it is \emph{primitive} if $A_1(a)=\F a$. The
algebra is a \emph{primitive $\mathcal F$-axial algebra} if it is generated by
primitive $\mathcal F$-axes. A \emph{Frobenius form} is a non-zero bilinear form
$(\cdot,\cdot)$ on~$A$ with $(uv,w)=(u,vw)$ for all $u,v,w\in A$.

Shpectorov asks~\cite[Problem~3.3]{gs}:
\begin{quote}
Does there exist an axial algebra with a Frobenius form such that the full
automorphism group of the algebra does not preserve the form?
\end{quote}
The survey adds that the example should be non-trivial: an algebra $A=B\oplus B$
with a form that is non-zero only on the first summand does not count, because
the automorphism exchanging the summands obviously moves the form. It is
especially interested in algebras whose space of Frobenius forms is
$1$-dimensional; then each automorphism $g$ multiplies the form by a scalar
$\alpha_g$, and for $g\mapsto\alpha_g$ to be non-trivial the algebra cannot have
an identity element. In most known examples the space of Frobenius forms is
$1$-dimensional and the full automorphism group preserves the form~\cite{gs}.
Searches of the literature on 12 and 13 November 2026 (the second covering the
arXiv listings announced on 13 November 2026) found no answer to the question,
and neither version of~\cite{gs} marks it as answered; the coverage of the
searches is described in the repository~\cite{repo}.

Let $\F=\F_{11}$ and let $A$ be the $4$-dimensional commutative algebra with
basis $a,b,c,d$ and product
\begin{gather*}
a^2=a,\quad b^2=b,\quad c^2=c,\quad d^2=d,\\
ab=c+d,\quad ac=-b-d,\quad ad=b+c,\quad bc=a+d,\quad bd=-a-c,\quad cd=a+b.
\end{gather*}
Let $\mathcal F=\{1,5,7,10\}\subseteq\F_{11}$ (note $10=-1$) with the fusion law
\begin{center}
\begin{tabular}{c|cccc}
$\ast$ & $1$ & $5$ & $7$ & $10$\\\hline
$1$ & $1$ & $5$ & $7$ & $10$\\
$5$ & $5$ & $1,5,7$ & $5,7$ & $10$\\
$7$ & $7$ & $5,7$ & $1,5,7$ & $10$\\
$10$ & $10$ & $10$ & $10$ & $1,5,7$
\end{tabular}
\end{center}
(the entry in row $\lambda$ and column $\mu$ lists $\lambda\ast\mu$). It is
finite and symmetric, and it is $C_2$-graded with even part $\{1,5,7\}$ and odd
part $\{10\}$ in the sense of~\cite[Section~2]{gs}. Let $\beta$ be the bilinear
form with $\beta(a,a)=\beta(c,c)=1$, $\beta(b,b)=\beta(d,d)=-1$ and
$\beta(x,y)=0$ for distinct basis vectors $x,y$, and let $g$ be the linear map
that exchanges $a$ with $b$ and $c$ with~$d$.

\begin{theorem}\label{thm:main}
\begin{enumerate}
\item[(a)] $A$ is a primitive $\mathcal F$-axial algebra, generated by the
primitive axes $a,b,c,d$.
\item[(b)] $\beta$ is a non-degenerate symmetric Frobenius form, and every
bilinear form $\gamma$ on $A$ with $\gamma(uv,w)=\gamma(u,vw)$ for all $u,v,w$
is a scalar multiple of~$\beta$.
\item[(c)] $g$ is an automorphism of $A$ of order~$2$, and
$\beta(gu,gv)=-\beta(u,v)$ for all $u,v\in A$.
\item[(d)] $A$ is simple and has no identity element; in particular $A$ is not
the direct sum of two non-zero ideals.
\end{enumerate}
So the full automorphism group of $A$ does not preserve its Frobenius form, the
space of Frobenius forms of $A$ (together with~$0$) is $1$-dimensional, and the
character $g\mapsto\alpha_g$ of~\cite{gs} takes the value $-1$ at~$g$.
\end{theorem}

The mechanism is visible in the table: $\beta$ gives the axes $a$ and $c$ length
$1$ and the axes $b$ and $d$ length $-1$, and $g$ exchanges the two pairs. We
have not determined the full automorphism group of~$A$; one element that moves
the form is enough.

\paragraph{Scope.}
The question asks only for an axial algebra, and the definitions
of~\cite{gs} allow any field and any fusion law; our example uses the field
$\F_{11}$ and a fusion law chosen for it. The law is not of Jordan type
$\{1,0,\eta\}$ or Monster type $\{1,0,\alpha,\beta\}$: it does not contain~$0$,
and the multiplication map of each of the four axes is invertible. For
algebras of Jordan type, \cite{gs} recalls a result of Hall, Segev and
Shpectorov: such an algebra has a unique Frobenius form in which every
primitive axis has length~$1$. Every automorphism preserves this form, so an
example of Jordan type would need a space of Frobenius forms of dimension at
least~$2$. We do not know whether there is an example of Monster type, or one
over a field of characteristic~$0$. The axes all have $1$-dimensional eigenspaces, so the fusion rules are a
finite computation; the law above lists, for each pair of eigenvalues, exactly
the eigenvalues that occur in the products.

\begin{remark}
The journal version of~\cite{gs} adds, after the question, a remark of
J.~McInroy on the Miyamoto group of an algebra with a $C_2$-graded fusion law
and the Frobenius form. In our example the Miyamoto involutions
(the maps that are $1$ on the even part $A_1(x)+A_5(x)+A_7(x)$ and $-1$ on the
odd part $A_{10}(x)$ of an axis~$x$) are $\tau_a=\tau_c$, which exchanges $b$
and~$d$, and $\tau_b=\tau_d$, which exchanges $a$ and~$c$; both preserve
$\beta$, they generate a group of order~$4$, and $g$ is not in it. This remark
was checked by computer and is not formalized.
\end{remark}

\section{Proof of Theorem~\ref{thm:main}}\label{sec:proof}

All coefficients are in $\F_{11}$. Coordinates are with respect to the basis
$a,b,c,d$.

\paragraph{(a) The axes.}
The four basis vectors are idempotents and form a basis, so they generate~$A$.
For each basis vector $x$, the following vectors are eigenvectors of $L_x$ for
the eigenvalues $1,5,7,10$:
\begin{center}
\begin{tabular}{ccccc}
\toprule
axis & $1$ & $5$ & $7$ & $10$\\\midrule
$a$ & $(1,0,0,0)$ & $(0,1,7,1)$ & $(0,1,5,1)$ & $(0,-1,0,1)$\\
$b$ & $(0,1,0,0)$ & $(8,0,8,1)$ & $(9,0,9,1)$ & $(-1,0,1,0)$\\
$c$ & $(0,0,1,0)$ & $(7,1,0,1)$ & $(5,1,0,1)$ & $(0,-1,0,1)$\\
$d$ & $(0,0,0,1)$ & $(1,7,1,0)$ & $(1,5,1,0)$ & $(-1,0,1,0)$\\
\bottomrule
\end{tabular}
\end{center}
This is checked by multiplying out. For each axis the four vectors have distinct
eigenvalues, so they form a basis; hence $L_x$ is diagonalizable with
eigenvalues $1,5,7,10$, each eigenspace is $1$-dimensional, and
$A_1(x)=\F x$. Write $v_\lambda$ for the listed $\lambda$-eigenvector of a
fixed axis $x$, so $v_1=x$ and $v_1v_\lambda=\lambda v_\lambda$. The other
products, written in the basis $(v_1,v_5,v_7,v_{10})$, are
\begin{center}
\begin{tabular}{ccc}
\toprule
product & axes $a,c,d$ & axis $b$\\\midrule
$v_5v_5$ & $(4,8,7,0)$ & $(3,9,7,0)$\\
$v_5v_7$ & $(0,6,7,0)$ & $(0,10,1,0)$\\
$v_5v_{10}$ & $(0,0,0,5)$ & $(0,0,0,7)$\\
$v_7v_7$ & $(7,6,5,0)$ & $(6,3,1,0)$\\
$v_7v_{10}$ & $(0,0,0,7)$ & $(0,0,0,8)$\\
$v_{10}v_{10}$ & $(2,4,8,0)$ & $(2,6,7,0)$\\
\bottomrule
\end{tabular}
\end{center}
(again by multiplying out). Since the eigenspaces are $1$-dimensional, this
gives $A_\lambda(x)A_\mu(x)\subseteq A_{\lambda\ast\mu}(x)$ for all
$\lambda,\mu\in\mathcal F$, so each basis vector is a primitive
$\mathcal F$-axis.

\paragraph{(b) The forms.}
Clearly $\beta$ is symmetric and non-degenerate. Let $C(x,y,z)=\beta(xy,z)$ for
basis vectors $x,y,z$. From the product, $C(a,a,a)=C(c,c,c)=1$,
$C(b,b,b)=C(d,d,d)=-1$, $C(x,y,z)=0$ when exactly two of $x,y,z$ are equal,
and for three distinct basis vectors
\[
C(a,b,c)=1,\quad C(a,b,d)=-1,\quad C(a,c,d)=1,\quad C(b,c,d)=-1,
\]
the same for every order of the three arguments. So $C$ is symmetric in its
three arguments, and $\beta(xy,z)=C(x,y,z)=C(y,z,x)=\beta(x,yz)$; by linearity
$\beta$ is a Frobenius form.

Now let $\gamma$ be any bilinear form with $\gamma(uv,w)=\gamma(u,vw)$; we do
not assume that $\gamma$ is symmetric. Fix a basis vector $x$ and let $y$ be
another one. Then $\gamma(x,y)=\gamma(x^2,y)=\gamma(x,xy)$, and $xy$ is plus or
minus the sum of the two basis vectors other than $x$ and~$y$. For $x=a$ this
gives, with $(p,q,r)=(\gamma(a,b),\gamma(a,c),\gamma(a,d))$,
\[
p=q+r,\qquad q=-p-r,\qquad r=p+q,
\]
which forces $2p=2q=2r=0$, so $p=q=r=0$. The same happens for $x=b,c,d$ (only
the signs change), so $\gamma(x,y)=0$ for all distinct basis vectors. Next,
$\gamma(ab,c)=\gamma(a,bc)$ gives $\gamma(d,c)+\gamma(c,c)=\gamma(a,a)+\gamma(a,d)$,
that is $\gamma(c,c)=\gamma(a,a)$; in the same way $\gamma(ab,d)=\gamma(a,bd)$
gives $\gamma(d,d)=-\gamma(a,a)$, and $\gamma(ac,b)=\gamma(a,cb)$ gives
$\gamma(b,b)=-\gamma(a,a)$. Hence $\gamma=\gamma(a,a)\,\beta$.

\paragraph{(c) The automorphism.}
The map $g$ is linear and $g^2=1$. It permutes the idempotents $a,b,c,d$, and on
the six products of distinct basis vectors
\begin{gather*}
g(ab)=d+c=ba,\quad g(ac)=-a-c=bd,\quad g(ad)=a+d=bc,\\
g(bc)=b+c=ad,\quad g(bd)=-b-d=ac,\quad g(cd)=b+a=dc,
\end{gather*}
so $g(uv)=g(u)g(v)$ for all $u,v$ by bilinearity. Since $g$ exchanges a basis
vector of length $1$ with one of length $-1$, $\beta(gu,gv)=-\beta(u,v)$.

\paragraph{(d) Simplicity and the identity.}
Fix a basis vector~$x$. By~(a), $L_x$ is diagonalizable with eigenvalues
$1,5,7,10$, so $(L_x-1)(L_x-5)(L_x-7)(L_x-10)=0$, which in $\F_{11}$ reads
$L_x^4-L_x^3+L_x^2+L_x=2$. Hence $L_x$ is invertible. The map
$P_x=3L_x^3+3L_x+6$ is $1$ on the eigenvalue~$1$ and $0$ on $5$, $7$ and~$10$,
so it is the projection onto $\F x$ along the other eigenspaces. In the table
of~(a), the eigenvectors of $x$ for $5$, $7$, $10$ have coordinate~$0$ at~$x$;
therefore $P_x(u)=u_x\,x$, where $u_x$ is the coordinate of $u$ at~$x$. Now
let $I$ be a non-zero ideal and $u\in I$ non-zero, say $u_x\ne0$. Since $I$ is
stable under $L_x$, it is stable under $P_x$, so $x\in I$, and then
$A=L_x(A)\subseteq I$. Since also $A^2\ne0$, $A$ is simple; in particular it
is not a direct sum of two non-zero ideals. If $z$ were an identity element, comparing the
$b$-coordinate of $zb=b$, the $d$-coordinate of $zd=d$ and the $c$-coordinate
of $za=a$ would give $z_b=1$, $z_d=1$ and $z_b+z_d=0$, so $2=0$ in $\F_{11}$, a
contradiction. \qed

\section{Formal verification}\label{sec:lean}

Theorem~\ref{thm:main} is formally verified in Lean~4 (version 4.34.1) with
Mathlib~\cite{mathlib} (version v4.34.1, commit \texttt{d13f23b7}). The files
are in version~\repoversion{} of the repository~\cite{repo}, directory
\texttt{lean/\allowbreak Research/\allowbreak TaskG/}.

\paragraph{The statement.}
The statement file \texttt{Challenge.lean} contains only definitions and the
statement, no proofs. It reuses the definitions of fusion law, axis, primitive
axis, primitive axial algebra, ideal, Frobenius form and the product given by
structure constants from the file \texttt{lean/\allowbreak Research/\allowbreak
AxialMSZ/\allowbreak Challenge.lean} of an earlier paper in the same
repository, which was reviewed and frozen then and is unchanged. The main
statement, in ASCII notation, is
{\footnotesize
\begin{verbatim}
def WitnessConditions (F : FusionLaw K) (mu : V ->l V ->l V) (X : Set V)
    (beta : V ->l V ->l K) (g : V =l V) (c : K) : Prop :=
  F.IsFiniteSymmetric /\ X.Finite /\ X.Nonempty /\
  IsPrimitiveAxialAlgebra F mu X /\ DirectlyIndecomposable mu /\
  beta != 0 /\ (forall u v, beta u v = beta v u) /\ IsFrobeniusForm mu beta /\
  IsAlgebraAutomorphism mu g /\ c != 0 /\ c != 1 /\
  (forall u v, beta (g u) (g v) = c * beta u v) /\
  ((forall u, g (g u) = u) \/ (forall u, g (g (g u)) = u))
def ExampleOver (K : Type*) [Field K] : Prop :=
  exists n T F X beta g c, 2 <= n /\ n <= 8 /\
    WitnessConditions F (structProduct T) X beta g c    -- V = Fin n -> K
def ShpectorovAnswer : Prop :=
  ExampleOver Rat \/ ExampleOver (ZMod 5) \/ ExampleOver (ZMod 7) \/ ExampleOver (ZMod 11)
\end{verbatim}}
\noindent In this notation, bilinear products have the type
\texttt{V ->l V ->l V}, invertible linear maps of $V$ have the type
\texttt{V =l V}, \texttt{IsAlgebraAutomorphism mu g} says $g(uv)=g(u)g(v)$, and
\texttt{DirectlyIndecomposable mu} says that whenever $A=I\oplus J$ for ideals
$I,J$, one of them is~$0$; this excludes every decomposition into two non-zero
ideals, not only $B\oplus B$. \texttt{IsPrimitiveAxialAlgebra F mu X} says
that the product is commutative, that every element of $X$ is a primitive axis
and that $X$ generates the algebra. \texttt{IsFiniteSymmetric} asks that the fusion
law be finite and symmetric, with values in~$\mathcal F$; a primitive axis is
a non-zero idempotent $a$ with $A=A_{\mathcal F}(a)$, the fusion rules, and
$A_1(a)=\F a$, as in Section~1. The bounds on $n$, the four fields and the
order of $g$ only narrow the search and make the statement stronger. The
stronger statement \texttt{StrongShpectorovAnswer} adds to
\texttt{WitnessConditions} that every bilinear form $\gamma$ (not assumed
symmetric) with $\gamma(uv,w)=\gamma(u,vw)$ is a scalar multiple of $\beta$,
and that the algebra has no identity element.

\paragraph{Review and freezing.}
An independent Codex agent, which wrote none of the files, reviewed the
statement against the question and the definitions of~\cite{gs} and approved
it, and the file was frozen on 12 November 2026 at 16:27~UTC, before any proof
file was written. The first compilation failed because Lean needs the facts that
$5$, $7$ and $11$ are prime to treat \texttt{ZMod 5}, \texttt{ZMod 7} and
\texttt{ZMod 11} as fields; three lines providing them (each proved by
\texttt{decide +kernel}) and a comment were added, nothing else changed (by
then the proof files had been drafted; the change concerns only these
instances), the same reviewer
checked the difference, and the file was frozen again at 16:41~UTC with the
SHA-256 hash
\begin{center}\small
\texttt{05448636ff2c4fefdc731111eaf9659228123084ecc24525cfde978a37709bb8}
\end{center}
of \texttt{Challenge.lean}. The files were not changed when they were moved to
the repository.

\paragraph{The proof.}
\texttt{Data.lean} defines $A$ (as $\F_{11}^4$ with the structure constants
above), $\beta$, $g$, the fusion law and the eigenvectors. The parts of
Theorem~\ref{thm:main} are the following theorems, in the namespace
\texttt{TaskG.Proof}, proved along the lines of Section~\ref{sec:proof}:
\begingroup\raggedright
\begin{itemize}
\item[(a)] \texttt{primitive\_basis\_axis}, \texttt{primitive\_axial} and
\texttt{fusion\_finite\_symmetric} in \texttt{Axes.lean}, and
\texttt{fusion\_grading} in \texttt{Main.lean};
\item[(b)] \texttt{beta\_symmetric}, \texttt{beta\_frobenius} and
\texttt{beta\_form\_space} in \texttt{Forms.lean}, and
\texttt{beta\_nondegenerate} in \texttt{Main.lean};
\item[(c)] \texttt{g\_mul}, \texttt{g\_involutive} and \texttt{g\_scales} in
\texttt{Automorphism.lean};
\item[(d)] \texttt{simple}, \texttt{directlyIndecomposable} and
\texttt{hasNoIdentity} in \texttt{Simplicity.lean}.
\end{itemize}
\endgroup
\texttt{Main.lean} combines them into proofs of \texttt{ShpectorovAnswer} and
\texttt{StrongShpectorovAnswer} with the field \texttt{ZMod 11} and $c=-1$.
The frozen statements only say that some example exists; that it is the
algebra of Theorem~\ref{thm:main} is seen from the definitions in
\texttt{Data.lean}, which \texttt{Check.lean} prints. Simplicity, the
non-degeneracy of~$\beta$ and the $C_2$-grading are proved about this algebra
in the proof files and are not part of the frozen statements. The finite
computations use \texttt{decide +kernel} and \texttt{ring}. The audit file
\texttt{Check.lean} restates both frozen statements and prints their axioms,
and the proof files print the axioms of the theorems above: each list is
exactly \texttt{propext}, \texttt{Classical.choice} and \texttt{Quot.sound}. No file
uses \texttt{sorry}, \texttt{native\_decide}, \texttt{bv\_decide} or an
\texttt{axiom} declaration.

\paragraph{Checks.}
On 12 November 2026 (17:25--17:32~UTC) ten files (the eight files of
\texttt{TaskG}, the statement file of \texttt{AxialMSZ} that they use, and its
tests) were compiled
module by module, each by its own Lean process, in two new empty output directories on the author's
Windows computer, with the pinned toolchain and Mathlib, which were not
rebuilt; each module was then replayed with \texttt{leanchecker}. All $40$
commands ended with exit code~$0$, with no error, no warning and no
\texttt{sorry}, and the compiled \texttt{.olean} files of the two directories
are byte for byte identical. From 12 to 13 November 2026 the whole Lean project
of the repository, with these eight files, was rebuilt from a new empty directory
on the author's Windows computer, module by module, each by its own Lean
process, and not with \texttt{lake}: all $3019$ modules compiled and $3018$ passed
\texttt{leanchecker}; the remaining one, the root of the project, has only
imports and no constants. Some modules of another paper of the repository ran
out of memory and were run again in a second pass; no Lean check rejected a
module. The formalization has not been checked
with comparator or with an independent kernel.

\section*{How this paper was produced}
This work was carried out on 12 and 13 November 2026 (UTC) in a workflow directed by
the author, using OpenAI Codex with the model GPT-6.1 Sol and Anthropic Claude
Code with the model Claude Opus~5.5. Claude Code's scouting agents proposed the
question, and Claude Code selected it and assigned it to Codex. Codex
agents found the example by a computer search over a family of
$4$-dimensional algebras over $\F_5$, $\F_7$ and $\F_{11}$, wrote the statement
file, the proofs and the Lean files, and ran the checks on the author's Windows
computer; an independent Codex agent (same model), which had written none of the
files, reviewed the statement, and another one gave the final review of the
mathematics, the Lean files and the records of the checks, recomputing the
example independently. Claude Code checked the example again by an independent
computer calculation, searched the literature, rebuilt the whole Lean project
of the repository and wrote this text; an independent Claude Code agent (model
Claude Opus~5.5), which had written none of the files, reviewed this paper and
the release before publication, and its corrections were applied. Throughout,
``we'' refers to this workflow; the reading and searching described in the paper
were done by the agents. The author is not a professional mathematician.
The author has read this paper and, using a side-by-side table, compared its
Lean statement with Theorem~\ref{thm:main}.
No human mathematician has reviewed this paper. The correctness of the result
rests on the Lean proof and on the agreement of the Lean statement with the
question of~\cite{gs}, which was checked by the reviewing agents; the written
proof is given for the reader. AI systems are not authors of this paper; the
author takes full responsibility for its content.

\medskip
\noindent Thanks are due to Ilya Gorshkov and Sergey Shpectorov for the
question and for collecting it with many others.

{\small
\begin{thebibliography}{99}
\bibitem{repo}
H.~Dong, \emph{ai4math-results}, repository,
\url{https://github.com/Horace-Maxwell/ai4math-results}, version~\repoversion{}
(2026); archived at Zenodo under the concept DOI
\url{https://doi.org/10.5281/zenodo.22970254}.
\bibitem{gs}
I.~Gorshkov and S.~Shpectorov, Axial algebras: questions and conjectures,
\emph{Commun.\ Math.} \textbf{35} (2027), no.~2, Paper no.~5,
\url{https://doi.org/10.46298/cm.19581}; arXiv:2606.30048v1, 2026. We use the
numbering of arXiv version~1: the question, attributed to S.~Shpectorov, is
Problem~3.3 there (in Section~3.1) and Problem~2.3 in the journal version,
with the same wording.
\bibitem{mathlib}
The mathlib Community, The Lean mathematical library, in \emph{Proc.\ 9th ACM
SIGPLAN Int.\ Conf.\ Certified Programs and Proofs} (CPP '20), ACM, 2020,
\url{https://doi.org/10.1145/3372885.3373824}.
\end{thebibliography}}
\end{document}
